\chapter{Conclusão}
\label{cap:conclusao}

Ao longo deste trabalho, foi documentada a jornada de ideação, planejamento e desenvolvimento do FaunaMar para auxiliar em um desafio real no campo do monitoramento ambiental. Este capítulo final visa consolidar os resultados alcançados, avaliar as limitações e levantar futuras implementações do sistema.

O projeto se deu a partir da percepção de um potencial de melhora no processo de coleta e gerenciamento de dados de monitoramento da fauna costeira, realizado pelo CECLIMAR. Um processo que era manual, fragmentado e suscetível a inconsistências, que demandava um esforço operacional alto. O objetivo principal, portanto, foi o desenvolvimento do aplicativo móvel FaunaMar, um sistema que, além de automatizar, padronizar e otimizar o ciclo de vida dos dados, busca também aumentar a participação da comunidade no processo de monitoramento ambiental a partir da ciência cidadã. Aliando assim a aplicação da tecnologia e inovação com a conservação da biodiversidade.

É possível afirmar que o objetivo geral foi alcançado. O principal resultado deste trabalho é uma aplicação funcional, robusta e pronta para uso, na qual a arquitetura foi detalhada no capítulo anterior. A automação da coleta e armazenamento foi implementada através de formulários de registro que persistem os dados diretamente no \textit{backend} ou localmente em modo \textit{offline}. A padronização dos registros foi garantida por meio de validações de formulário, campos de seleção baseados em dados pré-catalogados e uma estrutura de banco de dados NoSQL bem definida no \textit{Firestore}.

A promoção da participação da comunidade foi um dos pilares de ideação e construção da aplicação, que oferece uma interface amigável e intuitiva, um perfil com gamificação e um canal direto de contribuição para a pesquisa científica. Alguns termos técnicos que poderiam causar confusão foram evitados e textos de auxílio foram adicionados em algumas funcionalidades.

A facilitação da análise de dados foi implementada no "Painel de Registros", que oferece uma \textit{dashboard} com estatísticas, um mapa interativo, gráficos e uma funcionalidade de exportação de dados que respeita a privacidade dos usuários e pode ser utilizada por pesquisadores ou por cientistas cidadãos entusiastas.

É fundamental ressaltar que a qualidade e a maturidade do produto final estão diretamente ligados à metodologia de desenvolvimento adotada. A escolha pelo modelo iterativo incremental, gerenciado com Kanban, permitiu que o projeto evoluísse de forma adaptativa. Esse modelo, aliado à entrega contínua de 17 versões de teste para 7 usuários em dispositivos distintos, gerou um grande volume de \textit{feedbacks} que se transformaram em demandas. 42 dos 43 \textit{bugs} identificados e 9 das 19 sugestões de melhorias/débitos técnicos foram concluídos, além de refinamentos de usabilidade e da adição de funcionalidades não previstas no escopo inicial. Aliado a isso, o uso de um versionamento semântico facilitou o rastreamento das alterações e a comunicação com os testadores. Essa combinação entre o desenvolvimento iterativo, gerenciamento visual, testes contínuos e versionamento permitiu que o aplicativo \textit{FaunaMar}, se transformasse em uma solução mais alinhada às necessidades reais dos usuários.

Apesar dos avanços, o projeto apresenta limitações. A versão atual foi validada apenas para dispositivos Android, e a versão para iOS, embora viável com Flutter, não pôde ser testada por falta dos dispositivos necessários. A validação em iOS é importante, pois o sistema operacional tem particularidades que podem impactar a usabilidade e a performance do aplicativo. Sendo um sistema operacional amplamente utilizado, é importante garantir o funcionamento do FaunaMar em iOS para alcançar um público mais amplo.

Além disso, a escalabilidade da aplicação é um aspecto a ser acompanhado, pois a infraestrutura de \textit{backend} foi implementada utilizando o plano \textit{Spark}, a modalidade gratuita do \textit{Firebase}. Embora esse plano tenha se mostrado suficiente durante o lançamento inicial, ele possui limites de operações que irão afetar o desempenho em cenários de uso intensivo, especialmente com um grande volume de usuários simultâneos. Para garantir a escalabilidade da aplicação, a migração para o plano \textit{Blaze} é uma necessidade futura, permitindo que o projeto suporte um crescimento significativo sem comprometer a performance.

Como trabalhos futuros, o passo mais imediato é a preparação e publicação oficial do aplicativo na Google Play Store, para que a solução seja efetivamente disponibilizada para a comunidade de cientistas cidadãos. Em paralelo a isso, diversas sugestões de melhorias foram levantadas (conforme descritas na Seção~\ref{sec:gerencia-projeto}), mas não puderam ser implementadas no escopo do trabalho. Entre elas, a possibilidade de uma seção para "Fauna Local" que atualmente funciona com um redirecionamento externo, o que poderia ser substituído por uma visualização dinâmica dos dados da coleção \texttt{animals} diretamente no aplicativo após algumas alterações na coleção. Outras melhorias que se destacam são a adição de notificações \textit{push}, melhorias na lógica de navegação, melhor padronização de exceções e erros.

O \textit{FaunaMar} representa uma iniciativa de como a tecnologia pode ser aplicada de forma prática no auxílio à conservação da biodiversidade, ao mesmo tempo em que instiga a sociedade na produção de conhecimento científico. Dessa forma, o projeto não apenas otimiza o trabalho dos pesquisadores do CECLIMAR, como também contribui para a formação de uma comunidade mais engajada em iniciativas de cunho científico e ambiental.